\documentclass[10pt,a4paper]{amsart}
\usepackage[margin=19mm]{geometry}
\usepackage{amsmath,amssymb,mathtools}
\usepackage[scr=rsfs,cal=euler]{mathalfa}
\usepackage{xfrac}
\usepackage[unicode,hidelinks,bookmarks=false]{hyperref}
\newcommand{\ie}{i.e.\ }
\newcommand{\cf}{cf.\ }
\newcommand{\resp}{respectively\ }
\newcommand{\Subsection}{\subsection}
\providecommand{\nobreakdash}{\nobreak\hbox{-}}
\setlength{\emergencystretch}{2em}
\multlinegap0pt
\usepackage{bm}
\usepackage{bbm}
\usepackage{dsfont}
\usepackage{xspace}
\usepackage{cancel}
\usepackage{mathtools}
\usepackage{amsthm}
\makeatletter
\@ifundefined{Theorem}{\newtheorem{Theorem}{Theorem}}{}
\@ifundefined{Lemma}{\newtheorem{Lemma}[Theorem]{Lemma}}{}
\makeatother
\newcommand{\RR}{{{\rm I} \kern -.15em {\rm R} }}
\renewcommand{\geq}{\geqslant}
\renewcommand{\ge}{\geqslant}
\renewcommand{\leq}{\leqslant}
\renewcommand{\le}{\leqslant}
\title[Energy decay for semilinear evolution equations with memory ]{Energy decay for semilinear evolution equations with memory and time-dependent time delay feedback}
\author{Elisa Continelli, Cristina Pignotti}
\date{Source: arXiv:2404.17706v1 (CC BY 4.0)}
\begin{document}
\maketitle

\noindent\emph{Source and licence.} Energy decay for semilinear evolution equations with memory and time-dependent time delay feedback. Version arXiv:2404.17706v1 (CC BY 4.0), distributed under the Creative Commons Attribution 4.0 International licence, as recorded in the arXiv licence metadata for this exact version and shown on the abstract page https://arxiv.org/abs/2404.17706v1. Full rights notice: licenses/arxiv-2404.17706-CC-BY-4.0.txt.\par\smallskip

\noindent\textit{Editorial scope.} Counted unit: the Theorem whose source label is \texttt{teorema2.5} in the article (source0.tex lines 859--875, source label \texttt{teorema2.5}), delivered together with the article's own complete original proof of it (source0.tex lines 876--1081, one \texttt{proof} environment of 206 lines). No mathematics is written, repaired or re-derived: statement and proof are the article's own text line for line. Required context retained uncounted: modello (lines 199--210); K (lines 218--220); stima\_h (lines 281--284); modelloDafermos (lines 294--307); ipotesi2 (lines 337--340); forma\_astratta2 (lines 373--379); lemma1senzabound (lines 437--440); stimaE (lines 524--539); disuguaglianza energia (lines 531--534); Cbar (lines 536--538); Lemma 2 (lines 626--649); stima E dal basso (lines 636--643); assumptionPsi (lines 659--667); StabilityAbstract (lines 747--759); boundedsolution (lines 749--751); stimaesponenziale (lines 754--757). Every definition, display and label of those lines is retained unchanged. Omitted from this package and not redistributed: 1--198, 211--217, 221--280, 285--293, 308--336, 341--372, 380--436, 441--523, 535--535, 539--625, 644--658, 668--746, 752--753, 758--858, 1082--1536. Nothing inside a retained excerpt is omitted or altered: every retained body is delivered exactly as the article's lines are written, so no omission receipt of any kind accompanies this package. Citations: the retained excerpts carry no citation command at all, so no bibliography entry of the article is transcribed and no reference is redistributed. Reference adaptation: none was needed and none was made; every reference inside the delivered unit resolves inside it, so no unresolved reference or citation is produced. Printed slips kept literal: no printed word, symbol, hypothesis or number of the counted statement or of its proof was altered. Layout and class: the article's own class is \texttt{article} with packages including amsmath, amssymb, amsthm, graphicx, amscd, caption, amsfonts, mathrsfs, multicol, color, babel, fontenc, textcomp, inputenc; the wrapper is the lane's pinned \texttt{amsart} editorial wrapper at 10pt on A4, which restates the theorem-like environments the retained text needs and adds the article's own macro definitions that the retained text needs (\texttt{\textbackslash RR}, \texttt{\textbackslash ge}, \texttt{\textbackslash geq}, \texttt{\textbackslash le}, \texttt{\textbackslash leq}), with extra wrapper packages (bm, bbm, dsfont, xspace, cancel, mathtools). The delivered numbering is the unit's own. Assets: no retained span contains a figure environment, a graphics-inclusion command, a PDF-page inclusion, a table or a TikZ picture; no image file is copied into this package, the package declares no asset dependency and no image file is redistributed with it. Licence: version arXiv:2404.17706v1 of this article is distributed under the Creative Commons Attribution 4.0 International licence, as recorded in the arXiv licence metadata for this exact version and shown on the abstract page https://arxiv.org/abs/2404.17706v1. A full rights notice is delivered at licenses/arxiv-2404.17706-CC-BY-4.0.txt. Characters: every retained excerpt is pure ASCII, so the delivered engine, XeLaTeX with its default Latin Modern text font, reports no missing character.
\begin{equation}
\label{modello}
\begin{array}{l}
\displaystyle{u_{tt}(t)+Au(t)- \int_0^{+\infty} \beta (s)Au(t-s) ds+k(t)BB^*u_t(t-\tau (t))=\nabla \psi( u(t)), }\\
\hspace{12 cm}\ t\in (0,+\infty),\\
%\hspace{3 cm}
\displaystyle{
 u(t)=u_0(t), \quad t\in (-\infty, 0],}\\
%\hspace{3 cm}\\
\displaystyle{ u_t(t)=g(t), \quad t\in [-\bar\tau,0],}
\end{array}
\end{equation}

\begin{equation}\label{K}
\int_{t-\bar\tau}^t |k(s)| ds < K, \quad \forall t\geq 0.
\end{equation}

\begin{equation}
\label{stima_h}
||\nabla \psi (u)||_H\leq h(||A^{\frac 12} u||_H)||A^{\frac 12}u||_H,
\end{equation}

\begin{equation}
\label{modelloDafermos}
\begin{array}{l}
\displaystyle{u_{tt}(t)+(1-\tilde{\beta})Au(t)+\int_0^{+\infty} \beta (s)A\eta^t(s)ds+k(t)BB^*u_t(t-\tau(t))}\\
\displaystyle{\hspace{8 cm}
=\nabla \psi(u(t)),\ t\in(0,+\infty),}\\
\displaystyle{ \eta^t_t(s)=-\eta^t_s(s)+u_t(t),\ \quad t, s\in  (0, +\infty),}\\
%%\displaystyle{ u(x,t)=0 \quad \text{in} \quad \partial \Omega \times (0,+\infty),}\\
%%\displaystyle{ \eta^t(x,s)=0 \quad \text{in} \quad \partial\Omega \times (0,+\infty),}\\
\displaystyle{ u(0)=u_0(0),}\\
\displaystyle{u_t(t)= g(t), \quad t\in [-\bar\tau, 0],}\\
\displaystyle{ \eta^0(s)=\eta_0(s)=u_0(0)-u_0(-s)  \quad s\in (0,+\infty).}
\end{array}
\end{equation}

\begin{equation}
\label{ipotesi2}
b^2 Me^{\omega\bar\tau} \int_0^t  |k(s)| ds \leq \gamma +\omega 't, \quad \mbox{\rm for all}\ t\ge 0.
\end{equation}

\begin{equation}
\label{forma_astratta2}
\begin{array}{l}
\displaystyle{ U'(t)= \mathcal{A} U(t)-k(t)\mathcal{B}U(t-\tau(t))+F(U(t)),}\\
\displaystyle{ U(s)=\tilde g(s), \quad s\in [-\bar\tau, 0],}
\end{array}
\end{equation}

\begin{Theorem}
	\label{lemma1senzabound}
	Let us consider the system \eqref{forma_astratta2} with initial datum $\tilde g\in C([-\bar{\tau},0]; \mathcal{H}).$ Then, there exists a unique local solution $U(\cdot)$ defined on a time  interval $[0,\delta)$. 
\end{Theorem}

\begin{Lemma}\label{stimaE}
	Let $u:[0,T)\rightarrow \RR$ be a solution of \eqref{modello}.
	Assume that
	\begin{equation} \label{e(t)}
		E(t)\geq \frac{1}{4}||u_t(t)||_H^2,\quad \forall t\geq 0.
	\end{equation}
	Then,
	\begin{equation}
		\label{disuguaglianza energia}
		E(t)\leq \bar{C}(t)\mathcal{E}(0),\quad \forall t\geq 0,
	\end{equation}
	where
	\begin{equation}\label{Cbar}
		\bar{C}(t)=e^{3b^2\int_0^t|k(s)|  ds}.
	\end{equation}
\end{Lemma}

\begin{Lemma}
	\label{Lemma 2}
	Let $U(t)=(u(t),u_t(t),\eta^t)$ be a non-zero solution to \eqref{forma_astratta2} defined on an interval $[0, \delta),$ and let $T>\delta.$ Let $h$ be the strictly increasing function appearing in \eqref{stima_h}. 
	\begin{enumerate}
		\item If $h(||A^\frac{1}{2} u_0(0)||_H)<\frac{1-\tilde{\beta}}{2}$, then $E(0)>0$.
		\item Assume that $h(||A^\frac{1}{2} u_0(0)||_H)<\frac{1-\tilde{\beta}}{2}$ and that
		\begin{equation}\label{h1}
			h \left( \frac{2}{(1-\tilde{\beta})^\frac{1}{2}} C^{\frac{1}{2}}\mathcal{E}^\frac{1}{2}(0) \right) <\frac{1-\tilde{\beta}}{2},
		\end{equation}
		for some positive constant $C\geq \bar{C}(T)$, with $\bar{C}(\cdot)$ defined in  \eqref{Cbar}. Then
		\begin{equation}
			\label{stima E dal basso}
			\begin{array}{l}
				\vspace{0.3cm}\displaystyle{ E(t)>\frac{1}{4}||u_t(t)||_H^2+\frac{1-\tilde{\beta}}{4}||A^\frac{1}{2}u(t)||_H^2+\frac{1}{4}\int_{t-\bar\tau}^t |k(s)| \cdot ||B^*u_t(s)||_H^2 ds}\\
				\hspace{1.5 cm}
				\displaystyle{ +\frac{1}{4}\int_0^{+\infty} \beta(s) ||A^\frac{1}{2}\eta^t (s)||_H^2 ds,}
			\end{array}
		\end{equation}
		for all $t\in[0, \delta)$. In particular,
		\begin{equation}\label{J2}
			E(t)>  \frac 14 \Vert U(t)\Vert_{\mathcal H}^2, \quad \forall t\in [0, \delta).
		\end{equation}
	\end{enumerate}
\end{Lemma}

	\begin{equation}
		\label{assumptionPsi}
		\begin{array}{l}
			\displaystyle{|\psi(u)|\leq \int_0^1 |\langle \nabla \psi (su),u\rangle_H | ds} \\
			\hspace{1,15 cm}
			\displaystyle{\leq  ||A^\frac{1}{2}u||^2_H \int_0^1 h(s||A^\frac{1}{2}u||_H)sds}\\
			\hspace{1,15 cm}\displaystyle{\leq  h(||A^\frac{1}{2}u||_H)||A^\frac{1}{2}u||^2_H\int_{0}^{1}sds=\frac{1}{2}h(||A^\frac{1}{2}u||_H)||A^\frac{1}{2}u||^2_H,}
		\end{array}
	\end{equation}

\begin{Theorem}\label{StabilityAbstract}
Assume \eqref{ipotesi2}. Let $\tilde g\in C([-\bar\tau, 0];\mathcal{H})$ and let $U$ be a solution to \eqref{forma_astratta2} with the initial datum $\tilde{g}$, defined in a time interval $[0,T]$, $T>0$, that satisfies 
\begin{equation}\label{boundedsolution}
	\lVert U(t)\rVert_{\mathcal{H}}\leq C,\quad \forall t\in [0,T],
\end{equation}
for some $C>0$ such that $L(C)<\frac{\omega-\omega'}{M}$.
\\Then, $U$ satisfies the exponential decay estimate 
\begin{equation}
\label{stimaesponenziale}
\Vert  U(t)\Vert_{\mathcal H}\le Me^{\gamma}\left(||U_0||_{\mathcal H}+e^{\omega\bar{\tau}}Kb^2\max_{r\in[-\bar{\tau},0]}\{e^{\omega r}||\tilde{g}(r)||_{\mathcal H}\}\right)e^{-(\omega-\omega'-ML(C))t},
\end{equation}
for all $t\in [0,T]$.
\end{Theorem}

\begin{Theorem}\label{teorema2.5}
	Assume \eqref{ipotesi2}. There exist $\rho>0$ and $C_{\rho}>0$, with $L(C_{\rho})<\frac{\omega-\omega'}{2M}$, for which, if $\tilde g=(u_0(0),g,\eta_0)$ is such that
	\begin{equation}\label{smallness}
		\begin{array}{l}
			\displaystyle{||u_1||^2_{H}+(1-\tilde{\beta})||A^{\frac{1}{2}}u_0(0)||^2_{H}+\int_{-\bar\tau}^0|k(s)|\cdot||B^*g(s)||^2_{H}}\\
			\displaystyle{\hspace{5cm}+\int_{0}^{+\infty}\beta(s)||A^{\frac{1}{2}}\eta_0(s)||^2_{H}ds<\rho^2}
\end{array}
\end{equation}
and
\begin{equation}\label{smallness_bis}
			\displaystyle{\max_{s\in [-\bar{\tau},0]}||g(s)||_{H}<\rho,}
	\end{equation}
then the problem \eqref{modelloDafermos} with  initial datum $\tilde{g}$ has a unique global solution $U\in C([0,+\infty)),\mathcal{H})$ that satisfies
\begin{equation}\label{boundedsol2}
	\lVert U(t)\rVert_{\mathcal{H}}\leq C_{\rho},\quad \forall t\geq 0.
\end{equation} 
\end{Theorem}

\begin{proof}
	Let us fix a time $T$ sufficiently large, $T\ge \bar{\tau},$ such that
	\begin{equation}\label{CT}
		C_T:=4M^2e^{2\gamma}\max\left\{(1+Kb^2e^{\omega\bar{\tau}}),\, e^{\omega\bar{\tau}}\right\}\left(1+e^{2\omega\bar{\tau}}K^2b^4\right)e^{-(\omega-\omega')T}<1.
	\end{equation}
%%Note that, being $T$ large, we can assume, without loss of generality, that $T\geq %%\bar{\tau}$.
 Also, we set 
\begin{equation}\label{C*}
	C^*_T:=\sup\left\{e^{3b^2\int_{nT}^{(n+1)T}\lvert k(s)\rvert ds}: n\in \mathbb{N}\right\}.
\end{equation}
The assumption \eqref{K} ensures that  $C^*_T<+\infty.$
Note that  $C^*_T\ge \bar{C}(T),$ where $\bar{C}(T)$ is defined in \eqref{Cbar}.
Now, we pick $\rho>0$ in such a way that 
\begin{equation}\label{rho}
	\rho\leq \frac{(1-\tilde{\beta})^{\frac{1}{2}}}{2(C^*_T)^{\frac{1}{2}}}h^{-1}\left(\frac{1-\tilde{\beta}}{2}\right).
\end{equation} 
Let $(u_0(0),u_1,\eta_0)$ and $g$ be such that \eqref{smallness} and \eqref{smallness_bis}   holds true. Then, the initial datum $\tilde{g}$  satisfies the following smallness condition:
\begin{equation}\label{smallness2}
	\max_{s\in [-\bar{\tau},0]}||\tilde g(s)||_{\mathcal H}\leq \sqrt{2} \rho.
\end{equation}
Now, from Theorem \ref{lemma1senzabound} there exists a unique local solution $U(\cdot)$ to \eqref{forma_astratta2} with the initial datum $\tilde{g}(s)$, $s\in[-\bar{\tau},0]$ which is defined on a time interval $[0,\delta)$ and that satisfies the Duhamel's formula
\begin{equation}\label{primasol}
	U(t)=S(t)U_0+\int_{0}^{t}S(t-s)[-k(s)\mathcal{B}U(s-\tau(s))+F(U(s))]\,ds,
\end{equation} 
for all $t\in [0,\delta)$. Without loss of generality, we can suppose that $\delta\leq T$. Also, we can assume that $U$ is a non-zero solution to \eqref{forma_astratta2}, since, otherwise, \eqref{boundedsol2} is trivially satisfied. Thus, using \eqref{rho}, the assumption \eqref{smallness} on the initial data,
 and recalling that $h$ is a strictly increasing function, we get
$$h(||A^{\frac{1}{2}}u_0(0)||_{H})< h\left(\frac{\rho}{(1-\tilde{\beta})^{\frac{1}{2}}}\right)\leq h\left(\frac{1}{2(C_T^*)^{\frac{1}{2}}}h^{-1}\left(\frac{1-\tilde{\beta}}{2}\right)\right)\leq h\left(h^{-1}\left(\frac{1-\tilde{\beta}}{2}\right)\right)=\frac{1-\tilde{\beta}}{2},$$
were in the above inequality we used the fact that $2(C_T^*)^{\frac{1}{2}}\geq 1$. Hence, being $h(||A^{\frac{1}{2}}u_0(0)||_{H})< \frac{1-\tilde{\beta}}{2} $, from Lemma \ref{Lemma 2} it follows that $E(0)>0$. 
\\Also, using \eqref{assumptionPsi} and the fact that $h(||A^{\frac{1}{2}}u_0(0)||_{H})< \frac{1-\tilde{\beta}}{2} $, we can write
$$\begin{array}{l}
	\vspace{0.3cm}\displaystyle{E(0)=\frac{1}{2}||u_1||_H^2+\frac{1-\tilde{\beta}}{2}||A^\frac{1}{2}u_0(0)||_H^2-\psi(u_0(0))+\frac{1}{2}\int_{-\bar\tau}^0 |k(s)|\cdot ||B^*u_t(s)||^2_H ds}\\
	\vspace{0.3cm}\displaystyle{\hspace{2cm}+\frac{1}{2}\int_0^{+\infty} \beta(s)||A^\frac{1}{2}\eta_0(s)||_H^2 ds}\\
	\vspace{0.3cm}\displaystyle{\hspace{1cm}\leq\frac{1}{2}||u_1||_H^2+\frac{1-\tilde{\beta}}{2}||A^\frac{1}{2}u_0(0)||_H^2+\frac{1}{2}h(||A^\frac{1}{2}u_0(0)||_H)||A^\frac{1}{2}u_0(0)||^2_H}\\ \vspace{0.3cm}\displaystyle{\hspace{2cm}+\frac{1}{2}\int_{-\bar\tau}^0 |k(s)|\cdot ||B^*g(s)||^2_H ds+\frac{1}{2}\int_0^{+\infty} \beta(s)||A^\frac{1}{2}\eta_0(s)||_H^2 ds}\\
	\vspace{0.3cm}\displaystyle{\hspace{1cm}\leq\frac{1}{2}||u_1||_H^2+\frac{1-\tilde{\beta}}{2}||A^\frac{1}{2}u_0(0)||_H^2+\frac{1-\tilde{\beta}}{4}||A^\frac{1}{2}u_0(0)||^2_H}\\ \vspace{0.3cm}\displaystyle{\hspace{2cm}+\frac{1}{2}\int_{-\bar\tau}^0 |k(s)|\cdot ||B^*g(s)||^2_H ds+\frac{1}{2}\int_0^{+\infty} \beta(s)||A^\frac{1}{2}\eta_0(s)||_H^2 ds}\\
	\vspace{0.3cm}\displaystyle{\hspace{1cm}=\frac{1}{2}||u_1||_H^2+\frac{3}{4}(1-\tilde{\beta})||A^\frac{1}{2}u_0(0)||_H^2+\frac{1}{2}\int_{-\bar\tau}^0 |k(s)|\cdot ||B^*g(s)||^2_H ds}\\ \displaystyle{\hspace{2cm}+\frac{1}{2}\int_0^{+\infty} \beta(s)||A^\frac{1}{2}\eta_0(s)||_H^2 ds< \rho^2.}
\end{array} 
$$
The above inequality, together with \eqref{rho}, implies that
\begin{equation}\label{2}
	h\left(\frac{2}{(1-\tilde{\beta})^{\frac{1}{2}}}(C_T^*)^{\frac{1}{2}}E^{\frac{1}{2}}(0)\right)<h\left(\frac{2}{(1-\tilde{\beta})^{\frac{1}{2}}}(C_T^*)^{\frac{1}{2}}\rho\right)\leq h\left(h^{-1}\left(\frac{1-\tilde\beta}{2}\right)\right)=\frac{1-\tilde\beta}{2}.
\end{equation}
Furthermore, from \eqref{smallness_bis}, 
\begin{equation}\label{3}
	h\left(\frac{2}{(1-\tilde{\beta})^{\frac{1}{2}}}(C_T^*)^{\frac{1}{2}}\cdot\frac{1}{\sqrt{2}}\max_{s\in [-\bar{\tau},0]}||g(s)||_{H}\right)< h\left(\frac{\sqrt{2}}{(1-\tilde{\beta})^{\frac{1}{2}}}(C_T^*)^{\frac{1}{2}}\rho\right)\le \frac{1-\tilde\beta}{2}.
\end{equation}
Then, from \eqref{2}, \eqref{3} and by definition of $\mathcal{E}(0)$, we have that
\begin{equation}\label{4}
h\left(\frac{2}{(1-\tilde{\beta})^{\frac{1}{2}}}(C_T^*)^{\frac{1}{2}}\mathcal{E}^{\frac{1}{2}}(0)\right)<\frac{1-\tilde{\beta}}{2}.
\end{equation}
Since $C_T^*\geq \bar{C}(T)$, the above inequality \eqref{4} allows us to apply Lemma \ref{Lemma 2}, that ensures that \eqref{stima E dal basso} is satisfied for all $t\in [0,\delta)$. In particular,
$$E(t)\geq \frac{1}{4}||u_t(t)||_{H}^2,\quad \forall t\in [0,\delta).$$
Combining \eqref{disuguaglianza energia} and \eqref{stima E dal basso} with $\bar{C}(t)\leq \bar{C}(T)\leq C_T^*$, it turns out that
\begin{equation}\label{5}
	\begin{array}{l}
		\vspace{0.3cm}\displaystyle{\frac{1}{4}||u_t(t)||_H^2+\frac{1-\tilde{\beta}}{4}||A^\frac{1}{2}u(t)||_H^2+\frac{1}{4}\int_{t-\bar\tau}^t |k(s)| \cdot ||B^*u_t(s)||_H^2 ds
			+\frac{1}{4}\int_0^{+\infty} \beta(s) ||A^\frac{1}{2}\eta^t (s)||_H^2 ds}\\
		\displaystyle{\hspace{6cm}<E(t)\leq C_T^*\mathcal E(0),}
\end{array}
 \end{equation}
for all $t\in [0,\delta)$. Thus, since the solution $U$ is bounded from \eqref{5}, we can extend it in $t=\delta$ and in the whole interval $[0,T]$. Moreover, \eqref{5} holds for all $t\in [0,T]$.
\\Now, using \eqref{4} and \eqref{5} with $t=T$, we can write
\begin{equation}\label{t=T}
	h(||A^{\frac{1}{2}}u_0(T)||_{H})\leq h\left(\frac{2}{(1-\tilde{\beta})^{\frac{1}{2}}}(C_T^*)^{\frac{1}{2}}\mathcal{E}^{\frac{1}{2}}(0)\right)<\frac{1-\tilde{\beta}}{2}.
\end{equation}
Furthermore, from the smallness assumption \eqref{smallness} on the initial data and from \eqref{5}, it comes that
$$\frac{1}{4}||U(t)||^2_{\mathcal H}\leq E(t)\leq C_T^*\rho^2,$$
where here we have used the fact that ${\mathcal E}(0)<\rho^2$. Thus, 
\begin{equation}\label{boundedU}
	||U(t)||_{\mathcal H} \leq C_{\rho}:=2(C_T^*)^{\frac{1}{2}}\rho.
\end{equation}
Next, eventually choosing a smaller value of $\rho$, we can suppose that $L(C_\rho)<\frac{\omega-\omega'}{2M}$. We have so proved that there exist $\rho>0$, $C_\rho>0$ such that, whenever the initial data $(u_0(0),g,\eta_0)$ satisfy the smallness condition \eqref{smallness} and \eqref{smallness_bis}, then the system \eqref{forma_astratta2} with the initial data $U_0$, $\tilde{g}$, has a unique solution $U$ defined in the time interval $[0,T]$ such that $||U(t)||_{\mathcal H} \leq C_{\rho}$. 
\\As a consequence, from Theorem \ref{StabilityAbstract}, being $L(C_{\rho})<\frac{\omega-\omega'}{2M}$, the following estimate holds
$$\Vert  U(t)\Vert_{\mathcal H}\le Me^{\gamma}\left(||U_0||_{\mathcal H}+e^{\omega\bar{\tau}}Kb^2\max_{r\in[-\bar{\tau},0]}\{e^{\omega r}||\tilde{g}(r)||_{\mathcal H}\}\right)e^{-\frac{\omega-\omega'}{2}t},$$
for all $t\in [0,T]$. Thus, using \eqref{smallness2}, we can write
$$\Vert  U(t)\Vert_{\mathcal H}\le Me^{\gamma}\left(||U_0||_{\mathcal H}+e^{\omega\bar{\tau}}Kb^2 \sqrt{2} \rho\right)e^{-\frac{\omega-\omega'}{2}t},$$
for any $t\in [0,T]$. Then, 
$$\begin{array}{l}
	\vspace{0.3cm}\displaystyle{\Vert  U(t)\Vert^2_{\mathcal H}\le M^2e^{2\gamma}\left(||U_0||_{\mathcal H}+e^{\omega\bar{\tau}}Kb^2 \sqrt{2} \rho\right)^2e^{-(\omega-\omega')t}}\\
	\vspace{0.3cm}\displaystyle{\hspace{1.5cm}\leq2M^2e^{2\gamma}\left(||U_0||^2_{\mathcal H}+(e^{\omega\bar{\tau}}Kb^2 \sqrt{2} \rho )^2\right)e^{-(\omega-\omega')t}}\\
	\displaystyle{\hspace{1.5cm}=2M^2e^{2\gamma}\left(||U_0||^2_{\mathcal H}+e^{2\omega\bar{\tau}}K^2b^4 2\rho^2\right)e^{-(\omega-\omega')t},}
\end{array}$$
from which, taking into account \eqref{smallness2} with $U_0=\tilde{g}(0)$,
\begin{equation}\label{6}
	\Vert  U(t)\Vert^2_{\mathcal H}\le 4M^2e^{2\gamma}\rho^2\left(1+e^{2\omega\bar{\tau}}K^2b^4\right)e^{-(\omega-\omega')t},
\end{equation}
for all $t\in [0,T]$.
\\Also, for all $s\in [T-\bar{\tau},T]\subseteq [0,T]$, \eqref{6} yields
$$\begin{array}{l}
	\vspace{0.3cm}\displaystyle{||u_t(s)||_H^2\leq \Vert  U(s)\Vert^2_{\mathcal H}}\\
	\vspace{0.3cm}\displaystyle{\hspace{2.1cm}\leq 4M^2e^{2\gamma}\rho^2\left(1+e^{2\omega\bar{\tau}}K^2b^4\right)e^{-(\omega-\omega')s}}\\
	\vspace{0.3cm}\displaystyle{\hspace{2.1cm}\leq 4M^2e^{2\gamma}\rho^2\left(1+e^{2\omega\bar{\tau}}K^2b^4\right)e^{-(\omega-\omega')(T-\bar{\tau})}}\\
	\displaystyle{\hspace{2.1cm}\leq 4M^2e^{2\gamma}e^{\omega\bar{\tau}}\rho^2\left(1+e^{2\omega\bar{\tau}}K^2b^4\right)e^{-(\omega-\omega')T}.}
\end{array}$$
As a consequence, from \eqref{K}
$$\begin{array}{l}
	\vspace{0.3cm}\displaystyle{\int_{T-\bar{\tau}}^{T}|k(s)| \cdot ||B^*u_t(s)||_H^2 ds\leq 4M^2e^{2\gamma}e^{\omega\bar{\tau}}b^2\rho^2\left(1+e^{2\omega\bar{\tau}}K^2b^4\right)e^{-(\omega-\omega')T}\int_{T-\bar{\tau}}^{T}|k(s)|ds}\\
	\displaystyle{\hspace{3cm}\leq 4KM^2e^{2\gamma}e^{\omega\bar{\tau}}b^2\rho^2 \left(1+e^{2\omega\bar{\tau}}K^2b^4\right)e^{-(\omega-\omega')T}.}
\end{array}$$
This last fact together with \eqref{6} implies that
 \begin{equation}\label{7}
	\begin{array}{l}
		\vspace{0.3cm}\displaystyle{\Vert  U(T)\Vert^2_{\mathcal H}+\int_{T-\bar{\tau}}^{T}|k(s)| \cdot ||B^*u_t(s)||_H^2 ds}\\
		\displaystyle{\hspace{1cm}\leq 4M^2e^{2\gamma}\rho^2 \left(1+e^{2\omega\bar{\tau}}K^2b^4\right)\left(1+Kb^2e^{\omega\bar{\tau}}\right)e^{-(\omega-\omega')T}\leq C_T\rho^2<\rho^2.}
	\end{array}
\end{equation}
Moreover, \eqref{6} implies that
$$\max_{s\in [T-\bar{\tau},T]}||u_t(s)||_H^2\leq 4 M^2e^{2\gamma}e^{\omega\bar{\tau}}\rho^2 \left(1+e^{2\omega\bar{\tau}}K^2b^4\right)e^{-(\omega-\omega')T}\leq C_T\rho^2<\rho^2,$$
from which 
\begin{equation}\label{8}
	\max_{s\in [T-\bar{\tau},T]}||u_t(s)||_H<\rho.
\end{equation}
The conditions \eqref{7} and \eqref{8} allow us to apply the same arguments employed before on the interval $[T,2T]$. Namely, we consider  the initial value problem
\begin{equation}\label{modello2}
	\begin{array}{l}
		\displaystyle{V'(t)=\mathcal{A}V(t)-k(t)\mathcal{B}V(t-\tau(t))+F(V(t)),\quad t\in [T,2T],}\\
		\displaystyle{V(s)=U(s),\quad s\in [T-\bar{\tau},T],}
	\end{array}
\end{equation}
where $U(\cdot)$ is the solution to \eqref{forma_astratta2} in the interval $[0,T]$.
\\We define now the energy of the solution
\begin{equation}\label{energy}
	\begin{array}{l}
		\vspace{0.3cm}\displaystyle{\tilde{E}(t):=\tilde{E} (v(t))=\frac{1}{2}||v_t(t)||_H^2+\frac{1-\tilde{\beta}}{2}||A^{\frac 12}v(t)||^2_H-\psi(v)  }\\ \hspace{2 cm}
		\displaystyle{ +\frac{1}{2}\int_0^{+\infty} \beta(s) ||A^{\frac 1 2} \eta^t(s)||^2_H ds +\frac{1}{2}\int_{t-\bar\tau}^t |k(s)|\cdot ||B^*v_t(s)||_H^2 ds,}
	\end{array}
\end{equation}
where $\eta^t(s)=v(t)-v(t-s)$, and the functional
\begin{equation}\label{funzionale}
	\tilde{\mathcal{E}}(t):=\max\left\{\frac{1}{2}\max_{s\in [T-\bar{\tau},T]}\Vert u_t(s)\Vert_H^2,\,\,\max_{s\in[T,t]}\tilde{E}(s)\right\}.
\end{equation}
Let us note that $\tilde{E}(T)=E(T)$. 
\\Now, from Theorem \ref{lemma1senzabound}, the problem \eqref{modello2} with the initial datum $U(s)$, $s\in [T-\bar\tau,T]$, has a unique local solution $V(\cdot)$ defined on a time interval $[T,T+\delta)$ given by the Duhamel's formula
\begin{equation}\label{secondasol}
	V(t)=S(t-T)U(T)+\int_{0}^{t}S(t-T-s)[-k(s)\mathcal{B}V(s-\tau(s))+F(V(s))]\, ds,
\end{equation} 
for all $t\in [T,T+\delta)$. We can assume that $\delta\leq \bar{\tau}$ and that $V$ is a non-zero solution, since, otherwise, \eqref{boundedsol2} is obviously satisfied for all $t\geq T$.
\\First of all, the inequality \eqref{t=T} yields $\tilde{E}(T)>0$.
\\Let us note that, if $\tilde{E}(t)\geq \frac{1}{4}\lVert v_t(t)\rVert^2_{H}$, for all $t\in [T,T+\delta)$, arguing as in Lemma \ref{stimaE},
\begin{equation}\label{stimaEmodificata}
	\tilde{E}(t)\leq e^{3b^2\int_{T}^{t}\lvert k(s)\rvert ds}\,\tilde{\mathcal{E}}(T),\quad \forall t\in [T,T+\delta).
\end{equation}
Also, using the same arguments employed in Lemma \ref{Lemma 2} and observing that
$$C^*_T\ge e^{3b^2\int_T^{2T}\vert k(s)\vert ds },$$ 
if
\begin{equation}\label{h3}
	h\left(\frac{2}{(1-\tilde{\beta})^{\frac{1}{2}}}(C^*_T)^{\frac{1}{2}}\tilde{\mathcal{E}}^{\frac{1}{2}}(T)\right)<\frac{1-\tilde{\beta}}{2},
\end{equation}
 then 
\begin{equation}\label{stima E dal bassomodificata}
	\begin{array}{l}
		\vspace{0.3cm}\displaystyle{ \tilde{E}(t)>\frac{1}{4}||v_t(t)||_H^2+\frac{1-\tilde{\beta}}{4}||A^\frac{1}{2}v(t)||_H^2+\frac{1}{4}\int_{t-\bar\tau}^t |k(s)| \cdot ||B^*v_t(s)||_H^2 ds}\\
		\hspace{1.5 cm}
		\displaystyle{ +\frac{1}{4}\int_0^{+\infty} \beta(s) ||A^\frac{1}{2}\eta^t (s)||_H^2 ds,}
	\end{array}
\end{equation}
for all $t\in[T,T+ \delta)$. In particular,
\begin{equation}\label{J2modificata}
\tilde{E}(t)> \frac 14 \Vert V(t)\Vert_{\mathcal H}^2, \quad \forall t\in [T, T+\delta).
\end{equation}
Now, from \eqref{7}, we have $\tilde{E}(T)<\rho^2$. This together with \eqref{8} implies that
\begin{equation}\label{E<rho}
	\tilde{\mathcal{E}}(T)<\rho^2.
\end{equation}
Hence, using \eqref{rho} we get
 \begin{equation}\label{9}
	h\left(\frac{2}{(1-\tilde{\beta})^{\frac{1}{2}}}(C^*_T)^{\frac{1}{2}}\tilde{\mathcal{E}}(T)^{\frac{1}{2}}\right)<h\left(\frac{2}{(1-\tilde{\beta})^{\frac{1}{2}}}(C^*_T)^{\frac{1}{2}}\rho\right)<\frac{1-\tilde{\beta}}{2}.
\end{equation}
So, \eqref{h3} is satisfied. As a consequence, inequalities \eqref{stima E dal bassomodificata} and \eqref{J2modificata} hold for all $t\in [T,T+\delta)$. In particular, from \eqref{J2modificata} we get $$\tilde E(t)\geq \frac{1}{4}\lVert v_t(t)\rVert_H^2,\quad \forall t\in [T,T+\delta).$$
Therefore, also inequality \eqref{stimaEmodificata} is fulfilled. Combining \eqref{stimaEmodificata} and \eqref{stima E dal bassomodificata}, we finally get
\begin{equation}\label{solb}
	\begin{array}{l}
		\vspace{0.2cm}\displaystyle{\frac{1}{4}||v_t(t)||_H^2+\frac{1-\tilde{\beta}}{4}||A^\frac{1}{2}v(t)||_H^2+\frac{1}{4}\int_{t-\bar\tau}^t |k(s)| \cdot ||B^*v_t(s)||_H^2 ds+\frac{1}{4}\int_0^{+\infty} \beta(s) ||A^\frac{1}{2}\eta^t (s)||_H^2 ds}\\
		\displaystyle{\hspace{3cm}<\tilde{E}(t)\leq e^{3b^2\int_{T}^{2T}\lvert k(s)\rvert ds}\tilde {\mathcal{E}}(T)\leq C^*_T\tilde{\mathcal{E}}(T).}
	\end{array}
\end{equation}
Thus, the solution $V$ is bounded and we can extend it in $t=T+\delta$ and in the whole interval $[T,2T]$. Moreover, \eqref{solb} and \eqref{J2modificata} hold true for all $t\in [T,2T]$. So, taking into account of \eqref{E<rho}, it holds
\begin{equation}\label{boundedV}
	\lVert V(t)\rVert_{\mathcal{H}}\leq 2(C_T^*)^{\frac{1}{2}}\tilde{\mathcal{E}}^{\frac{1}{2}}(T)\leq 2(C_T^*)^{\frac{1}{2}}\rho=C_\rho.
\end{equation}
Putting together the two partial solutions \eqref{primasol} and \eqref{secondasol} to \eqref{forma_astratta2} obtained in the time intervals $[0,T]$ and $[T,2T]$, respectively, we get the existence of a unique solution $U\in C([0,2T];\mathcal{H})$ to \eqref{forma_astratta2} that is defined in the time interval $[0,2T]$ and that satisfies the Duhamel's formula \eqref{primasol}, for all $t\in [0,2T]$. Moreover, from \eqref{boundedU} and \eqref{boundedV}, the solution $U$ satisfies \eqref{boundedsolution} with $C=C_\rho$. Thus, since $L(C_\rho)<\frac{\omega-\omega'}{2M}$, the exponential decay estimate \eqref{stimaesponenziale} is satisfied by the solution $U$, i.e. $$\Vert  U(t)\Vert_{\mathcal H}\le Me^{\gamma}\left(||U_0||_{\mathcal H}+e^{\omega\bar{\tau}}Kb^2\max_{r\in[-\bar{\tau},0]}\{e^{\omega r}||\tilde{g}(r)||_{\mathcal H}\}\right)e^{-\frac{\omega-\omega'}{2}t},\quad\forall t\in [0,2T].$$
Again, we deduce that \eqref{6} holds for all $t\in [0,2T]$. Furthermore, for all $s\in [2T-\bar{\tau},2T],$
$$\begin{array}{l}
	\vspace{0.3cm}\displaystyle{||u_t(s)||_H^2 \leq ||U(s)||_{\mathcal{H}}^2\leq  4M^2e^{2\gamma}e^{\omega\bar{\tau}}\rho^2\left(1+e^{2\omega\bar{\tau}}K^2b^4\right)e^{-(\omega-\omega')2T}.}
\end{array}$$
As a consequence,
\begin{equation}\label{10}
\begin{array}{l}
		\vspace{0.3cm}\displaystyle{\Vert  U(2T)\Vert^2_{\mathcal H}+\int_{2T-\bar{\tau}}^{2T}|k(s)| \cdot ||B^*u_t(s)||_H^2 ds}\\
		\vspace{0.3cm}\displaystyle{\hspace{1cm}\leq 4 M^2e^{2\gamma}\rho^2(1+e^{2\omega \bar{\tau}}K^2b^4)e^{-(\omega-\omega')2T}}\\
		\vspace{0.3cm}\displaystyle{\hspace{2cm}+4 M^2e^{2\gamma}e^{\omega\bar{\tau}}Kb^2\rho^2 (1+e^{2\omega \bar{\tau}}K^2b^4)e^{-(\omega-\omega')2T}}\\
		\vspace{0.3cm}\displaystyle{\hspace{1cm}\leq 4 M^2e^{2\gamma}\rho^2 (1+e^{2\omega \bar{\tau}}K^2b^4)(1+Kb^2e^{\omega\bar{\tau}})e^{-(\omega-\omega')2T}}\\
		\vspace{0.3cm}\displaystyle{\hspace{1cm}\leq 4 M^2e^{2\gamma}\rho^2 (1+e^{2\omega \bar{\tau}}K^2b^4)(1+Kb^2e^{\omega\bar{\tau}})e^{-(\omega-\omega')T}}\\
		\displaystyle{\hspace{1cm}\leq C_T\rho^2<\rho^2.}
	\end{array}
\end{equation}
Also,
$$\begin{array}{l}
	\vspace{0.3cm}\displaystyle{\max_{s\in [2T-\bar{\tau},2T]}||u_t(s)||_H^2\leq 4M^2e^{2\gamma}e^{\omega\bar{\tau}}\rho^2 \left(1+e^{2\omega\bar{\tau}}K^2b^4\right)e^{-(\omega-\omega')2T}}\\
	\displaystyle{\hspace{2cm}\leq 4M^2e^{2\gamma}e^{\omega\bar{\tau}}\rho^2 \left(1+e^{2\omega\bar{\tau}}K^2b^4\right)e^{-(\omega-\omega')T}\leq C_T\rho^2<\rho^2,}
\end{array}$$
from which 
\begin{equation}\label{12}
	\max_{s\in [2T-\bar{\tau},2T]}||u_t(s)||_H<\rho.
\end{equation}
The conditions \eqref{10} and \eqref{12} allows us to apply the same arguments employed before on the interval $[2T,3T]$, obtaining the existence of a unique solution to \eqref{forma_astratta2} defined in the time interval $[0,3T]$ and that satisfies \eqref{boundedsol2} for all $t\in [0,3T]$. Iterating this procedure, we finally get the existence of a unique global solution $U\in C([0,+\infty);\mathcal{H})$ to \eqref{forma_astratta2} with the initial datum $\tilde{g}(s)$, $s\in [-\bar{\tau},0]$, that satisfies \eqref{boundedsol2}.
\end{proof}	

\end{document}
